\begin{proposition}
\label{proposition-finite-type-over-G-ring}
Let $R$ be a G-ring. If $R \to S$ is essentially of finite type
then $S$ is a G-ring.
\end{proposition}

\begin{proof}
Since being a G-ring is a property of the local rings it is clear
that a localization of a G-ring is a G-ring. Conversely, if every
localization at a prime is a G-ring, then the ring is a G-ring.
Thus it suffices to show that $S_\mathfrak q$ is a G-ring for every
finite type $R$-algebra $S$ and every prime $\mathfrak q$ of $S$.
Writing $S$ as a quotient of $R[x_1, \ldots, x_n]$ we see from
Lemma \ref{lemma-G-ring-goes-up-quasi-finite} that it suffices to prove
that $R[x_1, \ldots, x_n]$ is a G-ring. By induction on $n$ it
suffices to prove that $R[x]$ is a G-ring. Let $\mathfrak q \subset R[x]$
be a maximal ideal. By Lemma \ref{lemma-check-G-ring-maximal-ideals}
it suffices to show that
$$
R[x]_\mathfrak q \longrightarrow R[x]_\mathfrak q^\wedge
$$
is regular. If $\mathfrak q$ lies over $\mathfrak p \subset R$, then
we may replace $R$ by $R_\mathfrak p$. Hence we may assume that $R$
is a Noetherian local G-ring with maximal ideal $\mathfrak m$ and
that $\mathfrak q \subset R[x]$ lies over $\mathfrak m$. Note that
there is a unique prime $\mathfrak q' \subset R^\wedge[x]$
lying over $\mathfrak q$. Consider the diagram
$$
\xymatrix{
R[x]_\mathfrak q^\wedge \ar[r] &
(R^\wedge[x]_{\mathfrak q'})^\wedge \\
R[x]_\mathfrak q \ar[r] \ar[u] & R^\wedge[x]_{\mathfrak q'} \ar[u]
}
$$
Since $R$ is a G-ring the lower horizontal arrow is regular
(as a localization of a base change of the regular ring map
$R \to R^\wedge$). Suppose we can prove the right vertical arrow
is regular. Then it follows that the composition
$R[x]_\mathfrak q \to (R^\wedge[x]_{\mathfrak q'})^\wedge$
is regular, and hence the left vertical arrow is regular by
Lemma \ref{lemma-regular-permanence}.
Hence we see that we may assume $R$ is a Noetherian complete
local ring and $\mathfrak q$ a prime lying over the maximal
ideal of $R$.

\medskip\noindent
Let $R$ be a Noetherian complete local ring and let $\mathfrak q \subset R[x]$
be a maximal ideal lying over the maximal ideal of $R$. Let
$\mathfrak r \subset \mathfrak q$ be a prime ideal. We want to show that
$R[x]_\mathfrak q^\wedge \otimes_{R[x]} \kappa(\mathfrak r)$ is
a geometrically regular algebra over $\kappa(\mathfrak r)$.
Set $\mathfrak p = R \cap \mathfrak r$. Then we can replace $R$
by $R/\mathfrak p$ and $\mathfrak q$ and $\mathfrak r$ by their
images in $R/\mathfrak p[x]$, see
Lemma \ref{lemma-check-G-ring-easy}.
Hence we may assume that $R$ is a domain and that $\mathfrak r \cap R = (0)$.

\medskip\noindent
By  Algebra, Lemma
\ref{algebra-lemma-complete-local-Noetherian-domain-finite-over-regular}
we can find $R_0 \subset R$ which is regular and such that $R$ is
finite over $R_0$. Applying Lemma \ref{lemma-G-ring-goes-up-quasi-finite}
we see that it suffices to prove
$R[x]_\mathfrak q^\wedge \otimes_{R[x]} \kappa(\mathfrak r)$
is geometrically regular over $\kappa(\mathfrak r)$ when, in addition to the
above, $R$ is a regular complete local ring.

\medskip\noindent
Now $R$ is a regular complete local ring, we have
$\mathfrak r \subset \mathfrak q \subset R[x]$, we have
$(0) = R \cap \mathfrak r$ and $\mathfrak q$ is a maximal ideal
lying over the maximal ideal of $R$. Since $R$ is regular the
ring $R[x]$ is regular (Algebra, Lemma \ref{algebra-lemma-regular-goes-up}).
Hence the localization $R[x]_\mathfrak q$ is regular.
Hence the completions $R[x]_\mathfrak q^\wedge$ are regular, see
Lemma \ref{lemma-completion-regular}.
Hence the fibre $R[x]_{\mathfrak q}^\wedge \otimes_{R[x]} \kappa(\mathfrak r)$
is, as a localization of $R[x]_\mathfrak q^\wedge$, also regular.
Thus we are done if the characteristic of the fraction field of $R$ is $0$.

\medskip\noindent
If the characteristic of $R$ is positive, then $R = k[[x_1, \ldots, x_n]]$.
In this case we split the argument in two subcases:
\begin{enumerate}
\item The case $\mathfrak r = (0)$. The result is a direct consequence
of Lemma \ref{lemma-helper-G-ring}.
\item The case $\mathfrak r \not = (0)$. This is
Lemma \ref{lemma-another-helper-G-ring}.
\end{enumerate}
\end{proof}